\section{Conclusão}
\label{sec:conclusao}

A difusão da inteligência artificial na economia brasileira atinge primordialmente o segmento formal e escolarizado da força de trabalho, deparando-se com a informalidade contratual como amortecedor contra a automação corporativa. Após controlar para um vetor completo de características mincerianas (idade, idade ao quadrado, gênero e raça), cada ano adicional de escolaridade associa-se a 0{,}23 ponto a mais de exposição tecnológica em regressões no nível individual com erros-padrão agrupados por ocupação ($N = 227.629$). Em contrapartida, as ocupações que absorvem a maioria dos trabalhadores brasileiros (serviços domésticos, construção civil e agropecuária) operam com exposição quase nula.

A relação entre exposição e formalidade é substancialmente mediada pela escolaridade, e o gradiente é mais íngreme nas Unidades da Federação de setor formal mais profundo. O modelo teórico de designação de tarefas com adoção corporativa versus individual de IA explica esse conjunto de fatos: a informalidade atua como uma barreira protetiva não por vedar o acesso a ferramentas pontuais, mas por limitar a capacidade de extração de ganhos organizacionais de produtividade que justificam o custo de formalização.

Do ponto de vista de políticas públicas, os resultados reorientam o foco das iniciativas de capacitação e transição tecnológica. Como a informalidade amortece a exposição imediata, formalizar sem qualificação prévia pode expor trabalhadores a choques de automação desprovidos de ferramentas adaptativas, ao passo que capacitar sem o avanço da formalização restringe a captura de ganhos de escala corporativa. O público prioritário para programas de adaptação produtiva à IA não reside na base vulnerável do trabalho informal, mas no segmento médio-formalizado (técnicos, analistas e apoio administrativo), atualmente o mais sujeito à reorganização de rotinas informacionais. Programas de formação técnica e continuada --- em articulação com o Sistema S (SENAI/SENAC) e as redes públicas de educação profissional --- devem priorizar o desenvolvimento de competências complementares à IA para esse estrato intermediário.
